\documentclass[11pt,a4paper]{article}
\usepackage[utf8]{inputenc}
\usepackage[T1]{fontenc}
\usepackage[spanish]{babel}
\usepackage[margin=2.5cm]{geometry}
\usepackage{booktabs}
\usepackage{array}
\usepackage{tabularx}
\usepackage{listings}
\usepackage{xcolor}
\usepackage{enumitem}
\usepackage{parskip}
\usepackage{fancyhdr}
\usepackage{amsmath}
\usepackage{amssymb}
\usepackage{tikz}
\usetikzlibrary{arrows.meta, positioning}

\definecolor{codegray}{RGB}{245,245,245}
\definecolor{codeblue}{RGB}{0,60,180}
\definecolor{codegreen}{RGB}{0,120,0}

\lstdefinestyle{log}{
  basicstyle=\ttfamily\small,
  backgroundcolor=\color{codegray},
  frame=single,
  framerule=0.4pt,
  xleftmargin=8pt,
  numbers=none,
  breaklines=true
}

\pagestyle{fancy}
\fancyhf{}
\rhead{\textcolor{gray}{\small TD7: Ingeniería de Datos — UTDT}}
\lhead{\textcolor{gray}{\small Segundo Parcial Modelo 1}}
\rfoot{\thepage}
\renewcommand{\headrulewidth}{0.4pt}

\setlist[enumerate,1]{label=\textbf{\alph*)}}

\begin{document}

\begin{center}
  {\LARGE\bfseries TD7 — Ingeniería de Datos}\\[6pt]
  {\large Universidad Torcuato Di Tella}\\[4pt]
  {\Large\bfseries Segundo Parcial — Modelo 1}\\[4pt]
  {\normalsize Duración: 2 horas \quad|\quad Puntaje total: 100 puntos}
\end{center}
\noindent\rule{\linewidth}{0.8pt}
\vspace{2pt}

\noindent\textbf{Instrucciones:} Responda todos los ejercicios. Justifique cada respuesta. No está permitido el uso de material de consulta ni dispositivos electrónicos.

\vspace{6pt}
\noindent\rule{\linewidth}{0.4pt}

% ─────────────────────────────────────────────────────────────
\section*{Ejercicio 1 — Normalización \hfill\normalfont(35 puntos)}

Se tiene la siguiente relación sin normalizar:

\begin{center}
\textbf{INSCRIPCIÓN}(\texttt{legajo\_alumno, nombre\_alumno, email\_alumno, cod\_materia,}\\
\texttt{nombre\_materia, departamento\_materia, legajo\_docente, nombre\_docente, nota})
\end{center}

Las dependencias funcionales (DF) válidas son:

\begin{itemize}
  \item $\texttt{legajo\_alumno} \rightarrow \texttt{nombre\_alumno},\; \texttt{email\_alumno}$
  \item $\texttt{cod\_materia} \rightarrow \texttt{nombre\_materia},\; \texttt{departamento\_materia}$
  \item $\texttt{legajo\_docente} \rightarrow \texttt{nombre\_docente}$
  \item $\{\texttt{legajo\_alumno},\; \texttt{cod\_materia}\} \rightarrow \texttt{nota},\; \texttt{legajo\_docente}$
\end{itemize}

\begin{enumerate}
  \item (5 pts) Identificar la clave primaria de \textbf{INSCRIPCIÓN}. Justifique por qué es clave (es decir, por qué determina funcionalmente a todos los demás atributos).

  \vspace{3.5cm}

  \item (6 pts) Describir una anomalía de \textbf{inserción}, una de \textbf{modificación} y una de \textbf{eliminación} concretas que presenta esta relación.

  \vspace{4cm}

  \item (12 pts) La relación no está en 2FN. Explique por qué y llévela a \textbf{3FN}. Indique las nuevas relaciones con todos sus atributos, PK y FK.

  \vspace{7cm}

  \item (5 pts) ¿La descomposición obtenida en el punto anterior está en BCNF? Justifique revisando cada relación.

  \vspace{3.5cm}

  \item (7 pts) Verifique que la descomposición sea \textbf{sin pérdida} (\textit{lossless}). Aplique el algoritmo de la tabla de marcas o argumente a partir de los atributos compartidos entre las relaciones.

  \vspace{4cm}
\end{enumerate}

\newpage

% ─────────────────────────────────────────────────────────────
\section*{Ejercicio 2 — Concurrencia \hfill\normalfont(35 puntos)}

Se tienen tres transacciones con las siguientes operaciones:

\begin{center}
$T_1: R(X),\; R(Y),\; W(X)$ \qquad
$T_2: R(X),\; W(Y)$ \qquad
$T_3: R(Y),\; W(X)$
\end{center}

El siguiente schedule $S$ intercala sus operaciones:

\[
S:\; R_1(X),\; R_2(X),\; R_3(Y),\; W_2(Y),\; R_1(Y),\; W_1(X),\; W_3(X),\; C_1,\; C_2,\; C_3
\]

\begin{enumerate}
  \item (5 pts) Definir qué significa que dos operaciones estén en \textbf{conflicto}. ¿Por qué la serializabilidad por conflictos es una propiedad deseable?

  \vspace{3.5cm}

  \item (10 pts) Identificar \textbf{todos} los pares de operaciones en conflicto en $S$. Para cada par indique las operaciones, las transacciones y el ítem de datos involucrado.

  \vspace{5cm}

  \item (12 pts) Construir el \textbf{grafo de precedencias} de $S$ (nodos = transacciones, arcos según los conflictos). ¿$S$ es serializable por conflictos? Justifique. Si no lo es, identifique el ciclo.

  \vspace{5.5cm}

  \item (8 pts) ¿El schedule $S$ es \textbf{recuperable}? Para determinarlo, analice si hay alguna transacción que haya leído un ítem escrito por otra transacción sin que esta última haya hecho commit primero. Justifique con los pasos del schedule.

  \vspace{4.5cm}
\end{enumerate}

\newpage

% ─────────────────────────────────────────────────────────────
\section*{Ejercicio 3 — Recuperación y Almacenamiento \hfill\normalfont(30 puntos)}

\subsection*{Parte A — Recuperación con estrategia UNDO (12 puntos)}

El sistema utiliza la estrategia \textbf{UNDO} (actualización inmediata). Se tiene el siguiente log en disco al momento de la falla:

\begin{lstlisting}[style=log]
(BEGIN,  T1)
(WRITE,  T1, X, 10, 20)
(BEGIN,  T2)
(WRITE,  T2, Y, 5,  15)
(COMMIT, T1)
(BEGIN,  T3)
(WRITE,  T3, X, 20, 30)
         <-- FALLA DEL SISTEMA AQUI
\end{lstlisting}

\noindent\small\textit{Formato de WRITE: (WRITE, $T_i$, ítem, $v_{old}$, $v_{new}$).}

\normalsize
\begin{enumerate}
  \item (4 pts) ¿Qué transacciones deben deshacerse y por qué? ¿Cuál no debe tocarse?

  \vspace{2.5cm}

  \item (5 pts) Describir paso a paso el procedimiento de recuperación. ¿Cuáles son los valores finales de $X$ e $Y$ en disco luego de la recuperación?

  \vspace{3cm}

  \item (3 pts) ¿Por qué el algoritmo UNDO debe ser idempotente? ¿Qué pasaría si la recuperación se interrumpiera a mitad del proceso?

  \vspace{2.5cm}
\end{enumerate}

\subsection*{Parte B — Índices (8 puntos)}

\begin{enumerate}
  \setcounter{enumi}{3}
  \item (8 pts) Explicar la diferencia entre un índice \textbf{clustered} y uno \textbf{unclustered}. ¿Cuántos índices clustered puede tener una tabla? Dar un ejemplo concreto de cuándo es preferible cada tipo.

  \vspace{4.5cm}
\end{enumerate}

\subsection*{Parte C — Estimación de costos (10 puntos)}

Se tiene una relación $R$ con las siguientes estadísticas: $b_R = 200$ bloques, $f_R = 50$ tuplas/bloque (es decir, $n_R = 10\,000$ tuplas). Se realiza la consulta $\sigma_{\texttt{ciudad} = \texttt{'Rosario'}}(R)$. Se estima que 100 tuplas cumplen la condición. Asuma $t_T = 1\;\text{ms}$ y $t_s = 5\;\text{ms}$.

\begin{enumerate}
  \setcounter{enumi}{4}
  \item (3 pts) Calcular el costo de un \textbf{full scan}.

  \vspace{2.5cm}

  \item (4 pts) Calcular el costo usando un \textbf{índice B-tree clustered} sobre \texttt{ciudad}, con altura $h_i = 3$. Las 100 tuplas coincidentes ocupan 2 bloques contiguos.

  \vspace{2.5cm}

  \item (3 pts) Calcular el costo usando un \textbf{índice secundario} (unclustered) sobre \texttt{ciudad}, con altura $h_i = 3$, asumiendo que cada tupla coincidente puede estar en un bloque distinto.

  \vspace{2.5cm}
\end{enumerate}

\end{document}
